% Requires:
% \usetikzlibrary{calc}
\begin{tikzpicture}[
    scale=0.78,transform shape,
    block/.style={draw, fill=gray!20, rectangle, minimum width=1.6cm, minimum height=1.2cm, align=center, thick},
    wave arrow/.style={->, thick, blue, shorten >=1pt, shorten <=1pt},
    wavefront/.style={dashed,blue},
    separator/.style={dashed},
    structure/.style={draw, thick, line join=round, line cap=round},
    ac symbol/.style={thick},
    label text/.style={transform shape=false,align=center},
  ]

  % Region 1: Transmitter and Guided Waves

  % Transmitter Block
  \node[block] (transmitter) at (0,0) {};
  % AC Symbol
  \draw[ac symbol] (transmitter.center) circle (0.4cm);
  \draw[thick] (transmitter.center) ++ (-0.2,0) sin ++(+.1, .1) cos ++(+.1, -.1) sin ++(+.1, -.1) cos ++(+.1, +.1);
  \node[label text,below=0.15cm,inner sep=0pt] at (transmitter.south) {Transmitter};

  % Guided Waveguide (Structure)
  \draw[structure] (0.8, 0.4) -- (3, 0.4);
  \draw[structure] (0.8, -0.4) -- (3, -0.4);

  % Guided Waves (Arrows and Wavefronts)
  \foreach \x in {1.2, 1.8, 2.4} {
    \draw[wavefront] (\x, -0.4) -- (\x, 0.4);
    \draw[wave arrow] (\x-0.2, 0) -- (\x+0.2, 0);
  }
  % The last wavefront at the boundary
  \draw[wavefront] (3, -0.4) -- (3, 0.4);
  \draw[wave arrow] (2.8, 0) -- (3.2, 0);

  % Label for Guided Waves
  \node[label text,below=2.3cm,anchor=north] at (1.9, -0.2) {Guided\\1D\\waves};

  % Region 2:

  % Antenna Horn (Structure)
  \draw[structure] (3, 0.4) to[out=5, in=190] (6, 1.8);
  \draw[structure] (3, -0.4) to[out=-5, in=-190] (6, -1.8);
  % Antenna Aperture
  \draw[structure] (6, 0) ellipse (0.3cm and 1.8cm);
  % White fill to occlude the back of the aperture and the lines behind it
  \begin{scope}
    \clip (5.7, -1.9) rectangle (6, 1.9);
    \fill[white] (6, 0) ellipse (0.3cm and 1.8cm);
  \end{scope}
  % Redraw the front arc of the aperture
  \draw[structure] (6, -1.8) arc (270:90:0.3cm and 1.8cm);

  % Transition Waves (Bulging Wavefronts and Fanning Arrows)
  % Wavefront 1
  \draw[wavefront] (3.6, -0.5) to[out=65, in=295] (3.6, 0.5);
  \draw[wave arrow] (3.4, 0) -- (3.8, 0);

  % Wavefront 2
  \draw[wavefront] (4.4, -0.9) to[out=65, in=295] (4.4, 0.9);
  \draw[wave arrow] (4.2, 0) -- (4.6, 0);
  \draw[wave arrow] (4.2, 0.5) -- (4.6, 0.65);
  \draw[wave arrow] (4.2, -0.5) -- (4.6, -0.65);

  % Wavefront 3 (near aperture)
  \draw[wavefront] (5.4, -1.4) to[out=65, in=295] (5.4, 1.4);
  \draw[wave arrow] (5.1, 0) -- (5.6, 0);
  \draw[wave arrow] (5.1, 0.8) -- (5.5, 1.1);
  \draw[wave arrow] (5.1, -0.8) -- (5.5, -1.1);

  % Label for Transition Region
  \node[label text,below=2.3cm,anchor=north] at (4.5, -0.2) {Transition\\region\\(antenna)};

  % Region 3: Free Space

  % Center point for the free space arcs (virtual focal point)
  \coordinate (Source) at (3,0);

  % Free Space Waves (Concentric Arcs and Radial Arrows)
  % Wavefront 1
  \draw[wavefront] (Source) ++(30:4) arc (30:-30:4);
  \draw[wave arrow] ($(Source)+(0:3.8)$) -- ++(0:0.5);
  \draw[wave arrow] ($(Source)+(20:3.8)$) -- ++(20:0.5);
  \draw[wave arrow] ($(Source)+(-20:3.8)$) -- ++(-20:0.5);

  % Wavefront 2
  \draw[wavefront] (Source) ++(30:5.2) arc (30:-30:5.2);
  \draw[wave arrow] ($(Source)+(0:5.0)$) -- ++(0:0.5);
  \draw[wave arrow] ($(Source)+(20:5.0)$) -- ++(20:0.5);
  \draw[wave arrow] ($(Source)+(-20:5.0)$) -- ++(-20:0.5);

  % Wavefront 3
  \draw[wavefront] (Source) ++(30:6.4) arc (30:-30:6.4);
  \draw[wave arrow] ($(Source)+(0:6.2)$) -- ++(0:0.5);
  \draw[wave arrow] ($(Source)+(20:6.2)$) -- ++(20:0.5);
  \draw[wave arrow] ($(Source)+(-20:6.2)$) -- ++(-20:0.5);

  % Label for Free Space
  \node[label text,below=2.3cm,anchor=north] at (8.0, -0.2) {Free space\\3D\\waves};

  %--- Separators ---

  \draw[separator] (3, -2.5) -- (3, 2.5);
  \draw[separator] (6.3, -2.5) -- (6.3, 2.5); % Slightly after the antenna aperture

\end{tikzpicture}
